\clearpage
\VerdierPhysicalPageAnchor{0028}
\VerdierSemanticAnchor{intro-7-locally-compact-model}
局所コンパクト空間という特に簡単な場合が、得られる主張のモデルを与える。
$X,Y$ を二つの局所コンパクト空間（たとえば有限コホモロジー次元をもつ
もの）、$A$ を（Noether）環、$X_A,Y_A$ をそれぞれ $X,Y$ 上の
$A$-加群の層の圏、$\mathbf D(X_A),\mathbf D(Y_A)$ をその導来圏とし、
$f\colon X\longrightarrow Y$ を連続写像とする。この状況で次を定義できる：

\VerdierSemanticAnchor{intro-7-f-shriek}
\begin{enumerate}
\renewcommand{\labelenumi}{\alph{enumi})}
\item 関手
\[
f_!\colon X_A\longrightarrow Y_A,
\]
すなわち固有台付き順像と、それに対応する全導来関手
\VerdierSemanticAnchor{intro-7-rf-shriek}
\[
\mathrm Rf_!\colon \mathbf D(X_A)\longrightarrow \mathbf D(Y_A).
\]
（$f\colon X\longrightarrow Y$ が固有写像ならば、$f_!$ は通常の順像関手
である。）

\VerdierSemanticAnchor{intro-7-f-upper-shriek}
\item 関手
\[
f^!\colon \mathbf D^+(Y_A)\longrightarrow \mathbf D^+(X_A)
\]
（記号 $+$ は下に有界な複体に限ることを表す）。その厳密な定義はかなり
技術的だが、性質 c) によって一意同型を除いて特徴付けられる。

\VerdierSemanticAnchor{intro-7-duality-adjunction}
\item 二変数関手的な同型
\VerdierSemanticAnchor{intro-7-delta-f}
\[
\Delta_f(F,G)\colon
\operatorname{RHom}(\mathrm Rf_!F,G)
\longrightarrow
\operatorname{RHom}(F,f^!G)
\]
これは $\mathbf D(X_A)$ の任意の対象 $F$ と $\mathbf D^+(Y_A)$ の任意
の対象 $G$ に対して定義される（$\operatorname{RHom}$ は「準同型の複体」
関手の全導来関手を表す）。
\end{enumerate}

重要な現象は次のとおりである：

\VerdierSemanticAnchor{intro-7-duality-D1}
\noindent\textit{D1)} $f$ を変化させるとき、$f^!$ の形成は推移的である。

\VerdierSemanticAnchor{intro-7-duality-D2}
\noindent\textit{D2)} 同型 $\Delta_f$ は写像の合成と両立する。

\VerdierSemanticAnchor{intro-7-duality-D3}
\noindent\textit{D3)} $X$ が $Y$ 上相対次元 $n$ で滑らかなとき、すなわち
$X$ 上局所的に、$X$ が $Y$ のある開部分と位相空間 $\mathbf R^n$ との
積に $Y$ 上同型であるとき（例：$Y$ は一点、$X$ は位相多様体）、
$\mathbf D^+(Y_A)$ の任意の対象 $G$ に対して標準同型
\VerdierSemanticAnchor{intro-7-relative-orientation}
\[
f^!(G)\simeq f^*(G)\otimes\omega_{X/Y}[n]
\]
がある。ここで $f^*(G)$ は通常の逆像、$\omega_{X/Y}$ は階数 $1$ の
局所自由層、すなわち $Y$ 上の $X$ の相対向き付け層であり、
$\omega_{X/Y}[n]$ は複体 $\omega_{X/Y}$ の次数を $n$ だけシフトして
得られる複体である。
